% GENERATED FILE. Edit canonical lessons, then run npm run generate:books.
% canonical-source-hash: fnv1a64:1280a2c9c2d10f56
% canonical-lessons: TA-C270-kaikkuttai, TA-C270-totu2, TA-C270-cankili, TA-C270-cappatti, TA-C270-puri

\chapter{Handkerchief, Earring, Chain, Chapati, Puri}
\label{ch:ta-handkerchief-earring-chain-chapati-puri}
\hlchaptermodality{270}

\begin{chapteropening}
\textbf{By the end of this chapter:} \emph{I can say a handkerchief, an earring, a chain, a chapati and a puri.}

Complete the last lesson of chapter 270: I can say a handkerchief, an earring, a chain, a chapati and a puri.
\end{chapteropening}

\section[\texorpdfstring{\ta{கைக்குட்டை} (kaikkuṭṭai)}{கைக்குட்டை (kaikkuṭṭai)}]{\ta{கைக்குட்டை} (kaikkuṭṭai) — a handkerchief}

\label{lesson:TA-C270-kaikkuttai}



\begin{quote}
Before the new one: say the Tamil for simple, then the Tamil for comfortable.
\end{quote}

\subsection*{You'll want to know: \ta{கைக்குட்டை}}
\textbf{\ta{கைக்குட்டை}} — \emph{kaikkuṭṭai} — \textquotedblleft{}a handkerchief\textquotedblright{}.

\begin{grammarlens}[title={What you've built}]
One more word for this chapter.
\end{grammarlens}

\subsection*{Your turn}
\begin{itemize}
\raggedright
  \item \emph{Say it:} \emph{kaikkuṭṭai}
  \item \emph{Say it:} \emph{kaikkuṭṭai}, once more
  \item \emph{Say it:} say \emph{kaikkuṭṭai}
  \item \emph{Recall:} say the Tamil for simple, then the Tamil for comfortable, then say \emph{kaikkuṭṭai} again
\end{itemize}

\begin{tcolorbox}[breakable,colback=teal!4,colframe=teal!35!black,title={Before you move on}]
What does \ta{கைக்குட்டை} mean? (\textquotedblleft{}A handkerchief\textquotedblright{}.) Say it once more.
\end{tcolorbox}
\section[\texorpdfstring{\ta{தோடு} (tōṭu)}{தோடு (tōṭu)}]{\ta{தோடு} (tōṭu) — an earring}

\label{lesson:TA-C270-totu2}



\begin{quote}
Before the new one: say the Tamil for comfortable, then the Tamil for a handkerchief.
\end{quote}

\subsection*{You'll want to know: \ta{தோடு}}
\textbf{\ta{தோடு}} — \emph{tōṭu} — \textquotedblleft{}an earring\textquotedblright{}.

\begin{grammarlens}[title={What you've built}]
One more word for this chapter.
\end{grammarlens}

\subsection*{Your turn}
\begin{itemize}
\raggedright
  \item \emph{Say it:} \emph{tōṭu}
  \item \emph{Say it:} \emph{tōṭu}, once more
  \item \emph{Say it:} say \emph{tōṭu}
  \item \emph{Recall:} say the Tamil for comfortable, then the Tamil for a handkerchief, then say \emph{tōṭu} again
\end{itemize}

\begin{tcolorbox}[breakable,colback=teal!4,colframe=teal!35!black,title={Before you move on}]
What does \ta{தோடு} mean? (\textquotedblleft{}An earring\textquotedblright{}.) Say it once more.
\end{tcolorbox}
\section[\texorpdfstring{\ta{சங்கிலி} (caṅkili)}{சங்கிலி (caṅkili)}]{\ta{சங்கிலி} (caṅkili) — a chain, a necklace}

\label{lesson:TA-C270-cankili}



\begin{quote}
Before the new one: say the Tamil for a handkerchief, then the Tamil for an earring.
\end{quote}

\subsection*{You'll want to know: \ta{சங்கிலி}}
\textbf{\ta{சங்கிலி}} — \emph{caṅkili} — \textquotedblleft{}a chain, a necklace\textquotedblright{}.

\begin{grammarlens}[title={What you've built}]
One more word for this chapter.
\end{grammarlens}

\subsection*{Your turn}
\begin{itemize}
\raggedright
  \item \emph{Say it:} \emph{caṅkili}
  \item \emph{Say it:} \emph{caṅkili}, once more
  \item \emph{Say it:} say \emph{caṅkili}
  \item \emph{Recall:} say the Tamil for a handkerchief, then the Tamil for an earring, then say \emph{caṅkili} again
\end{itemize}

\begin{tcolorbox}[breakable,colback=teal!4,colframe=teal!35!black,title={Before you move on}]
What does \ta{சங்கிலி} mean? (\textquotedblleft{}A chain, a necklace\textquotedblright{}.) Say it once more.
\end{tcolorbox}
\section[\texorpdfstring{\ta{சப்பாத்தி} (cappātti)}{சப்பாத்தி (cappātti)}]{\ta{சப்பாத்தி} (cappātti) — a chapati}

\label{lesson:TA-C270-cappatti}



\begin{quote}
Before the new one: say the Tamil for an earring, then the Tamil for a chain.
\end{quote}

\subsection*{You'll want to know: \ta{சப்பாத்தி}}
\textbf{\ta{சப்பாத்தி}} — \emph{cappātti} — \textquotedblleft{}a chapati\textquotedblright{}.

\begin{grammarlens}[title={What you've built}]
One more word for this chapter.
\end{grammarlens}

\subsection*{Your turn}
\begin{itemize}
\raggedright
  \item \emph{Say it:} \emph{cappātti}
  \item \emph{Say it:} \emph{cappātti}, once more
  \item \emph{Say it:} say \emph{cappātti}
  \item \emph{Recall:} say the Tamil for an earring, then the Tamil for a chain, then say \emph{cappātti} again
\end{itemize}

\begin{tcolorbox}[breakable,colback=teal!4,colframe=teal!35!black,title={Before you move on}]
What does \ta{சப்பாத்தி} mean? (\textquotedblleft{}A chapati\textquotedblright{}.) Say it once more.
\end{tcolorbox}
\section[\texorpdfstring{\ta{பூரி} (pūri)}{பூரி (pūri)}]{\ta{பூரி} (pūri) — a puri, a puffed fried bread}

\label{lesson:TA-C270-puri}



\begin{quote}
Before the new one: say the Tamil for a chain, then the Tamil for a chapati.
\end{quote}

\subsection*{You'll want to know: \ta{பூரி}}
\textbf{\ta{பூரி}} — \emph{pūri} — \textquotedblleft{}a puri, a puffed fried bread\textquotedblright{}.

\begin{grammarlens}[title={What you've built}]
One more word for this chapter.
\end{grammarlens}

\subsection*{Your turn}
\begin{itemize}
\raggedright
  \item \emph{Say it:} \emph{pūri}
  \item \emph{Say it:} \emph{pūri}, once more
  \item \emph{Say it:} say \emph{pūri}
  \item \emph{Recall:} say the Tamil for a chain, then the Tamil for a chapati, then say \emph{pūri} again
\end{itemize}

\begin{tcolorbox}[breakable,colback=teal!4,colframe=teal!35!black,title={Before you move on}]
What does \ta{பூரி} mean? (\textquotedblleft{}A puri, a puffed fried bread\textquotedblright{}.) Say it once more.
\end{tcolorbox}
